\worldchapter{Einleitung}{introduction}
\index{Einleitung}
\label{ch:introduction}

\begin{multicols}{2}
Willkommen bei \textit{\trans{book_title}}!

\textit{\trans{book_title}} ist ein Rollenspiel. Diese Regeln und alle zugehörigen Materialien reichen aus, um Abenteuer in der Welt \trans{book_title} zu spielen. Du kannst diese Regeln frei verwenden, und mit Freunden Abenteuer in der Welt bestreiten. 

\wbif{phasesix}{Oder du erschaffst dir eine eigene Welt auf der Basis der in diesem System verfügbaren Epochen.}{} 

\begin{pquote}{Johann W. von Goethe, Faust}
Ihr wißt, auf unsren deutschen Bühnen\\    
Probiert ein jeder, was er mag;\\    
Drum schonet mir an diesem Tag\\    
Prospekte nicht und nicht Maschinen\\    
Gebraucht das groß' und kleine Himmelslicht,\\    
Die Sterne dürfet ihr verschwenden;\\    
An Wasser, Feuer, Felsenwänden,\\    
An Tier und Vögeln fehlt es nicht.\\    
So schreitet in dem engen Bretterhaus\\    
Den ganzen Kreis der Schöpfung aus\\    
Und wandelt mit bedächt'ger Schnelle\\    
Vom Himmel durch die Welt zur Hölle.    
\end{pquote}

Dieses Rollenspiel lässt sich komplett ohne elektronische Hilfsmittel spielen. Der Fokus des Systems liegt allerdings in unseren modernen Zeiten darauf, über eine Internetplattform wie Discord, Teamspeak oder Ähnlichem gespielt zu werden. Die zugehörige Plattform \href{\trans{book_website}}{\trans{book_website}} bietet alle möglichen Werkzeuge zur Verwaltung von Charakteren und Kampagnen. Zudem kann hier gewürfelt werden, Würfelergebnisse können optional an einen Discord Chat weitergegeben werden.

\section{Was ist Rollenspiel?}
\index{Was ist Rollenspiel?}

Um genau zu sein, reden wir hier von Pen and Paper Rollenspiel, also keinem Computerspiel. Traditionell spielt man ein Rollenspiel mit 2-4 \textit{Spielern} und einer \textit{Spielleiterin} am Tisch, wobei gegenüber einem Brettspiel das Spielbrett durch Charakterbögen und Würfel ersetzt werden. Das Rollenspiel erzählt immer eine Geschichte, die von allen Mitspielern getragen und weitergesponnen wird.

Die Spieler erschaffen für eine Rollenspielrunde Charaktere, also fiktive Figuren, welche auf Charakterbögen festgehalten werden. Der Charakterbogen enthält die Beschreibung des Charakters, seine Herkunft, seine Interessen und seine Fähigkeiten. Letzte sind in Zahlenwerten festgehalten, denn das Handeln im Rollenspiel erfordert sogenannte \textit{Proben} oder \textit{Würfe}, welche den Ausgang einer Aktion bestimmen.

\wbif{tirakan}{
\begin{pexample}[Beispiel]
Tom hat sich entschlossen, in einer Rollenspielgruppe mit Spielleiterin Julia mitzuspielen. Julia hat für ihre Runde ausgesucht, dass es ein Abenteuer im Königreich Asgoran des Jahres 322 geben wird.

Tom beschließt, als Charakter einen adeligen Ritter namens \textit{Jamie} zu erstellen. Er wählt die Charakterschablonen ``Ritter'' und ``Adelig'' und die Abstammungsschablone ``Asgoran'', und stellt mit diesen seinen Charakter fertig. Er addiert alle Werte der Schablonen und verzeichnet diese auf seinem Charakterbogen.

Durch den gewählten Hintergrund hat Jamie besonders hohe Anfangswerte in Fertigkeiten wie \textit{Nahkampf}, \textit{Ausdauer} und \textit{Geschick}. Bei \textit{Naturkunde} oder der \textit{Darbietung} ist sein Können jedoch allenfalls rudimentär.
\end{pexample}
}{
\begin{pexample}[Beispiel]
    Tom hat sich entschlossen, in einer Rollenspielgruppe mit Spielleiterin Julia mitzuspielen. Julia hat für ihre Runde ausgesucht, dass es ein Abenteuer\wbif{phasesix}{ auf der Erde des Jahres 1982}{} geben wird. Da es eine Geschichte im Stil der Retro-Science Fiction Mystery geben soll, gibt Julia die Epoche ``Der kalte Krieg und die 80er'' sowie die ``Horror Erweiterung'' vor.

Tom beschließt also, als Charakter den Journalisten Jamie mit Abitur zu erstellen. Er wählt die Charakterschablonen ``Journalist'' und ``Abitur'', und stellt mit diesen seinen Charakter fertig. Er addiert alle Werte der Schablonen und verzeichnet diese auf seinem Charakterbogen.

Durch den gewählten Hintergrund hat Jamie besonders hohe Anfangswerte in Fertigkeiten wie \textit{Untersuchen}, \textit{Kommunikation} und \textit{Politik}. Beim \textit{Schießen} oder der \textit{Akrobatik} ist sein Können jedoch allenfalls rudimentär.
\end{pexample}
}

Näheres zu der Erschaffung eines Charakters findest du in \cref{ch:create} und \cref{ch:rolls}.

Während die Spieler jeweils einen Charakter für das Spiel erstellen, bereitet die Spielleiterin eine Geschichte vor. Diese wird oft auch \textit{Abenteuer} oder \textit{Plot} genannt. Diese Geschichte ist nicht, wie in einem Roman, bis ins Letzte ausformuliert. Stattdessen ist es ein grobes Script, bestehend aus einem möglichen Ablauf, der Beschreibung von Orten, sowie sogenannten Nichtspielercharakteren (auch NSC oder NPC).

Sobald das Spiel beginnt, agieren alle in der Rolle ihrer Charaktere. Die Spielleiterin beschreibt Situationen mit Worten, Karten oder Zeichnungen. Die Spieler sprechen für ihre Charaktere in der ersten Person (``Ich schleiche die Treppe hinauf.''). Kommt es zu Handlungen der Charaktere mit ungewissem Ausgang, so wird auf Proben zurückgegriffen, und es wird gewürfelt.

\begin{pexample}[Beispiel]
Nachdem die Vorbereitungen abgeschlossen sind, trifft sich die Gruppe um Spielleiterin Julia. Dann geht es endlich los, und Julia beschreibt die erste Szene.
\end{pexample}

\wbif{tirakan}{
\begin{pexample}
Julia: ``Es ist der 2. Nebelmond des Jahres 322. Ihr befindet euch in einer Taverne in dem beschaulichen Städtchen Lindfield im Süden Asgorans. Es ist nun bereits spät am Abend, und vor der Tür hat der leichte Nieselregen dafür gesorgt, dass die Schneedecke der letzten Tage von einer dünnen Eisschicht bedeckt ist. Es wird kalt heute Nacht, und glatt. Die Taverne ist gut gefüllt. Mit einem Knarren öffnet sich die Eingangstür, und eine Wolke feinen Regens kommt in den Pub. Direkt gefolgt von einer Gestalt in einer viel zu kleinen Teerjacke.'' 
\end{pexample}
}{
\begin{pexample}
Julia: ``Es ist der 2. Januar \wbif{phasesix}{1982}{}. Ihr befindet euch in einem Pub in dem beschaulichen Städtchen Lindfield im Süden Englands. Es ist nun bereits spät am Abend, und vor der Tür hat der leichte Nieselregen dafür gesorgt, dass die Schneedecke der letzten Tage von einer dünnen Eisschicht bedeckt ist. Es wird kalt heute Nacht, und glatt. Der Pub ist gut gefüllt, und ihr hört Tainted Love aus einer Jukebox, während ihr auf ein neues Pint wartet. Mit einem Knarren öffnet sich die Eingangstür, und eine Wolke feinen Regens kommt in den Pub. Direkt gefolgt von einer Gestalt in einer viel zu kleinen, gelben Plastikregenjacke.''
\end{pexample}
}

\begin{pexample}
Dies ist also der Auftakt, und Tom entschließt sich, dass sein Charakter Jamie einmal einen Blick auf den Neuankömmling werfen möchte. Er sagt die Handlungen für Jamie an: 

Tom: ``Ich betrachte den Fremden da mal sehr genau, diese schlecht sitzende Jacke ist mir ja schon aufgefallen.'' 

Julia: ``Du erkennst, dass unter der Kapuze nasse, schwarze Haare in das Gesicht eines alten Mannes fallen. Wirf doch einmal \textit{Wahrnehmung}, um mehr zu erkennen.''
\end{pexample}

Es geht also um ein kooperatives Weiterentwickeln der Geschichte durch das Handeln der Charaktere. Die Spielleiterin hat einen Plan, wohin sich die Geschichte entwickeln könnte, welche Personen auftreten können, und was diese eigentlich für eine Motivation haben. Irgendein Geschehen entwickelt sich um die Charaktere der Spieler herum, und sie werden mit in diese Handlung hinein gezogen.

Wohin diese Geschichte führt, ist ungewiss. Es mag sein, dass etwas Schlimmes verhindert wird, oder dass ein Geheimnis aufgedeckt wird. Die Spielleiterin hat einen groben Plan, die Spieler bestimmen jedoch hauptsächlich den Fortgang.

\subsection{Es geht um das Erzählen} 
\index{Es geht um das Erzählen}

Denkt man an Computerrollenspiele, so ist die strategische Weiterentwicklung des Charakters der wichtigste Punkt. Er muss zukünftige Kämpfe bestehen können und möglichst optimale Werte für mögliche Herausforderungen haben. Im Pen and Paper Rollenspiel geht es um den Fortgang der Geschichte, um gemeinsame Erlebnisse und Erinnerungen. Das möglichst optimale Ausrichten auf ``starke'' Eigenschaften (sogenanntes \textit{Power-Gaming}) sollte hier nicht im Vordergrund stehen. Dadurch, dass die Geschichte immer gemeinsam weitergetragen wird, gibt es für alle Herausforderungen sehr flexible Lösungen.

\textbf{Ein Wort zum Power-Gaming}: Das Regelwerk von \trans{book_title} unterbindet bewusst nicht die Möglichkeit, einen relevanten Wert (z.B. \textit{Schießen}) in astronomische Höhen zu treiben. Es sollte in der Spielgruppe Einigkeit darüber bestehen, welchen Stil man beim Spiel haben möchte. Die Regeln ermöglichen solche Konstruktionen bewusst, um bei der Erschaffung von Charakteren und Abenteuern alle Freiheiten zu lassen.

Ebenso kommt hier die alte Rollenspielregel zum Tragen, dass das \textbf{Wort der Spielleiterin immer mehr wiegt als die Regeln}. Natürlich sollte es der Normalfall sein, dass die Regeln so angewendet werden, wie sie geschrieben sind, denn es geht für die Spieler auch um einen Rahmen, auf den sie sich verlassen können. Kommt es jedoch zu einer unklaren Regelsituation, so entscheidet die Ansage der Spielleiterin.

\subsection{Kampf im Rollenspiel} 
\index{Kampf im Rollenspiel}

Auch wenn der Fokus im Pen and Paper Rollenspiel weniger auf der Auseinandersetzung mit Fäusten oder Waffen liegt, spielt der Kampf doch eine wichtige Rolle. Nicht jede Situation lässt sich ohne Waffengewalt lösen. Schnell kommt es zu einer Schlägerei, oder die Charaktere planen, einen Händler zu überfallen.

Kampf im Rollenspiel wird anders behandelt als das freie Spiel. Die Zeit läuft in Runden ab, und man veranschaulicht die Situation auf einer Karte auf dem (eventuell virtuellen) Tisch. Spieler agieren nacheinander, die Spielleiterin steuert ihre NSC. Wunden zeigen an, wie gut es den Charakteren noch geht. Näheres zum Ablauf des Kampfs findest du in \cref{ch:combat}.

Im Spiel sollte sich beides die Waage halten. Es mag auch Abenteuer geben, die nur aus einem Kampf bestehen, \trans{book_title} ist jedoch keine realistische Kampfsimulation. Hier geht es darum, einen Konflikt möglichst unterhaltsam, cineastisch, und trotzdem spannend auszutragen.

Bei einem Kampf im \textit{\trans{book_title}} System sollten aufgrund der Besonderheiten (Reaktionen, Aktionen stehlen etc.) allerdings immer folgende Dinge eingehalten werden: 

\begin{itemize}
    \item Verwendet immer eine Karte. Eine Grundrisskarte der Situation sorgt dafür, dass es keine Missverständnisse bei der Positionierung gibt, egal wie kurz der Kampf ist. Eine Karte kann eine vorgefertigte, aufwändige Karte, aber auch ein schnell gezeichneter Grundriss sein. 
    \item Verwendet immer einen Massstab. Die Charaktere haben unterschiedliche Bewegungsreichweiten, und bei einigen ist es auch ihre Stärke, dass sie sich besonders weit bewegen können. Diesen Charakteren und Gegnern wird ohne Massstab ihre Stärke genommen.
    \item Benutzt einen Initiative-Tracker. Initiative bestimmt im Kampf die Reihenfolge der Beteiligten. Initiative-Tracker bedeutet, dass diese Reihenfolge für alle sichtbar aufgeschrieben ist. Es ist bei \trans{book_title} für die Spieler wichtig zu wissen, wann sie selbst wieder an der Reihe sind (denn dann verfallen u.A. aufgehobene eigene Handlungen).
\end{itemize}

\section{Besonderheiten} 
\index{Besonderheiten}

\trans{book_title} hat in einigen Bereichen andere Ansätze als andere Rollenspielsysteme. Teilweise wurden diese so gestaltet, um die Flexibilität bei den möglichen Szenarien zu erreichen. Auch legt das System einen großen Fokus darauf, möglichst schnell spielbar zu sein, und z.B. im Kampf heldenhafte Aktionen ausführen zu können.

Für Würfe und Proben werden normale, sechsseitige Würfel verwendet. Es werden Würfel in der Anzahl des jeweiligen Wertes geworfen. Jeder Würfel, der nach dem Wurf eine 5 oder höher zeigt, stellt einen Erfolg dar. In der Regel reicht ein einzelner Erfolg aus, um eine Probe als bestanden gelten zu lassen.

\subsection{Charakterschablonen} 
\index{Charakterschablonen}

Charaktere werden in \trans{book_title} nicht durch das Verteilen von Punkten auf Fertigkeiten, Attribute oder andere Werte erstellt oder gesteigert. Stattdessen werden \textit{Charakterschablonen} verwendet, die jeweils kleine Stationen in der Vergangenheit des Charakters darstellen.

Charakterschablonen sind unterteilt in \textit{Abstammung}, \textit{Beruf}, \textit{Bildung}, \textit{Charakter}, \textit{Talent} und \textit{Lebensumstände}. Während die Schablonen in den ersten beiden Kategorien viele Eigenschaften mit sich bringen (Ein Heiler bekommt Gewissenhaftigkeit, Erste Hilfe und Medizinkenntnisse), kann eine Schablone aus dem Bereich Talent zum Beispiel ``Guter Redner'' sein, welche Punkte auf \textit{Kommunikation} bringt.

Charakterschablonen werden für \textit{Reputation} erstanden, dies sind die Erfahrungspunkte, welche die Charaktere für bestandene Abenteuer erhalten.

\subsection{Besondere Möglichkeiten im Kampf} 
\index{Besondere Möglichkeiten im Kampf}

Der Kampf ist so gehalten, dass das Geschehen möglichst eindrucksvoll ist, die Mechanik aber gut von der Hand geht. Die in Rollenspielen übliche Reihenfolge der Teilnehmer ist hier auch zu finden, jedoch ist der ganze Ablauf ein wenig dynamischer. \textit{Reaktionen} sind in das Kampfsystem eingebaut, jede Aktion kann bis zum erneuten Begin der eigenen Kampfrunde aufgespart werden, und als Reaktion verwendet werden.

Zudem ist es möglich, durch das Ausgeben von \textit{Bonuswürfeln}, welche Charaktere durch Schablonen erhalten können, spontan eigene Aktionen im Kampf zu erzeugen oder durch \textit{Schicksalswürfel} sogar Gegnern Aktionen zu stehlen. Obwohl Schicksalswürfel sehr selten sind, kann dem Gegner so eventuell der tödliche Schlag gestohlen und zur eigenen Aktion gemacht werden. Was eventuell etwas unrealistisch klingt, sorgt so für einen sehr dynamischen Ablauf des Kampfes, es kann oft zu kinoreifen Szenen kommen, die auch die Spielleiterin nicht vorhersehen kann.

\subsection{Waffen} 
\index{Waffen}

Waffen sind in \trans{book_title} erweiterbar gestaltet. Es gibt eine Liste von Waffenmodifikationen. Unterschiedliche Munition wird auch als Waffenmodifikation abgebildet.

\wbif{phasesix}{
So bringt die Horrorerweiterung zum Beispiel Silbermunition mit, welche gut gegen Werwölfe wirkt. Durch den modularen Aufbau lässt diese sich sowohl im Mittelalter bei Bögen, als auch in der Moderne in einem Sturmgewehr nutzen.
}{}

Näheres zu den Kampfregeln und Waffenmodifikationen findest du in \cref{ch:combat}.

\wbif{phasesix}{
\subsection{Epochen und Erweiterungen} 
\index{Epochen und Erweiterungen}

\trans{book_title} ist als Regelwerk darauf ausgelegt, möglichst flexibel zu sein. Es kann für viele Szenarien verwendet werden, egal ob es sich um Fantasy, Science Fiction, Horror oder Geschichten in der ``echten'' Welt handelt. Dabei bietet es einen Grundstock an fertigen Waffen, Charakterschablonen, Gegenständen und Rüstungen, die in irdische Epochen eingeteilt sind. Zudem ist es natürlich möglich, dass eine Spielgruppe eigene Inhalte erschafft und verwendet.

Um zu erreichen, dass jedes Szenario möglich ist, unterscheidet \trans{book_title} drei Arten von Erweiterungen.

\subsubsection{Die Grundregeln} 
\index{Grundregeln}

Einige Elemente sind immer und überall gültig. Sie gelten unabhängig davon, welche Epoche oder Erweiterung gewählt ist. Viele Charakterschablonen wie ``Gewissenhaft'', ``Waffennarr'', ``Tratschtante'', aber auch Waffen und Gegenstände sind unabhängig von Epoche oder Erweiterung immer verfügbar.

\subsubsection{Epochen} 
\index{Epochen}

Epochen sind irdische Zeiträume, welche die Vorlage für alle Szenarien (auch Fantasy) sind. Sie geben einen technologischen Stand bei Waffen und Gegenständen vor und bestimmen, was für die Spieler verfügbar ist. Die irdische Geschichte ist in 8 Epochen eingeteilt.

\begin{itemize}
    \item \textbf{Die klassische Antike} - 800 v.Chr-600 n.Chr. 
    \item \textbf{Mittelalter, Vikinger und Kreuzzüge} - 500-1500 
    \item \textbf{Viktorianisches Zeitalter und der Wilde Westen} - 1700-1900 
    \item \textbf{Imperialismus und Weltkriege} - 1900-1950 
    \item \textbf{Der kalte Krieg und die 80er} - 1950-1990 
    \item \textbf{Moderne Zeit} - 1990 und danach 
%    \item \textbf{Nahe Zukunft} - eine dystopische nahe Zukunft
    \item \textbf{Science Fiction} - eine ferne Zukunft 
\end{itemize}

Die Inhalte der Epochen richten sich nach der irdischen Technologie der jeweiligen Zeit. Ein Abenteuer findet immer in einer der Epochen statt.

Zudem sind die Inhalte der Epochen sind möglichst frei von speziell irdischen Elementen gehalten, damit sie auch in einer eigenen Fantasywelt verwendet werden können. Natürlich hat zum Beispiel die Moderne Zeit bekannte moderne Waffen, und der Zweihänder ist auch eine irdische Erfindung. Es ist jedoch möglichst generisch gehalten, sodass es auch in ein Szenario passt, das nicht auf der Erde spielt.

\subsubsection{Erweiterungen} 
\index{Erweiterungen}

Zusätzlich zu den Epochen können bestimmte Erweiterungen gewählt werden, um einem Abenteuer etwa Magie oder das Wirken von Göttern hinzuzufügen. Diese Erweiterungen können von der Spielleiterin beliebig ausgewählt werden und sind optional.
\begin{itemize}
    \item \textbf{Magie} - fügt die magische Resource ``Arkana'' für die Charaktere hinzu und bringt Zauber und Artefakte mit sich.
    \item \textbf{Horror} - definiert Regeln für den Umgang mit Horrorelementen, Stress und Eigenarten 
    \item \textbf{Pantheon} - bietet Regeln für die Interaktion mit Göttern: Anrufungen, Gebete und Gunst 
    \item \textbf{Körpermodifikationen} - bringt einen Katalog von biomechanischen Elementen mit, die nach bestimmten Regeln in den Körper integriert werden können.
%    \item \textbf{Fahrzeuge und Dronen} - bietet spezielle Regeln für Fahrzeuge, Fahrzeugkampf und Aufbauten auf Fahrzeugen.
\end{itemize}

\subsection{Welten} 
\index{Welten}

Durch die Kombination von Epoche und Erweiterungen kann so also ein beliebiges Szenario geschaffen werden. Eine Cthulhu Geschichte im Wilden Westen ist ebenso möglich wie eine magische Welt in der Moderne. Eine klassische, eigene Fantasywelt könnte sich der Mittelalter-Epoche und der Erweiterung ``Magie'' bedienen.

Einige bestehende Welten fassen diese Kombination aus Epochen und Erweiterungen zusammen, und bringen zudem noch die Beschreibung einer ganzen Welt mit sich. Sie stehen als ganzes zur Verfügung und sind direkt benutzbar.

\subsubsection{Tirakans Reiche} 
\index{Tirakans Reiche}

Die Welt Tirakan ist eine komplette Fantasywelt, in der zu einem beliebigen Zeitpunkt in einer 1000-jährigen Geschichte gespielt werden kann. Eine ausgearbeitete Geschichte von Menschen, Elfen, Gnomen und vielen anderen Völkern erzählt den Kampf der Zivilisationen gegen Minotauren, Echsen und eine namenlose Finsternis.
\begin{itemize}
    \item \textbf{Epoche}: Mittelalter 
    \item \textbf{Erweiterungen}: Magie, Pantheon 
    \item \textbf{Weltbeschreibung}: \href{https://tirakans-reiche.de}{tirakans-reiche.de} 
\end{itemize}

\subsubsection{Die Abenteuer des Abteilung V des NEXUS} 
\index{Abenteuer des Abteilung V des NEXUS}

Die Geschichte der Abteilung V des NEXUS spielt in der Moderne. Es ist eine fiktive Geheimorganisation, die gegründet wurde um die Menschheit vor außerirdischen und paranormalen Gefahren zu schützen. Die Spieler spielen Agenten der Abteilung V des NEXUS, und erleben durch die Fähigkeit des Zeitreisens Abenteuer in allen möglichen Epochen und Welten.
\begin{itemize}
    \item \textbf{Epoche}: Moderne 
    \item \textbf{Erweiterungen}: Horror, Körpermodifikationen, Fahrzeuge 
    \item \textbf{Weltbeschreibung}: \href{https://phasesix.org/world/nexus}{phasesix.org} 
\end{itemize}
}{}
\end{multicols}
